\problema{Binary Tree Zigzag Level Order Traversal}{103}{src/02-arboles/103-zigzag-level-order.cpp}{M}

Dada la raíz de un árbol binario, queremos su recorrido por niveles, pero
alternando la dirección: el nivel 0 se lee de izquierda a derecha, el nivel 1
de derecha a izquierda, el nivel 2 otra vez de izquierda a derecha, y así
sucesivamente.

\paragraph{La idea, con un ejemplo de la vida real.} Primero, unas palabras.
Un \emph{nivel} del árbol es todo lo que está a la misma altura: la raíz es el
nivel $0$, sus hijos el nivel $1$, los nietos el nivel $2$, y así. Recorrer «por
niveles» es como leer un edificio piso por piso: primero toda la planta baja,
luego todo el primer piso, etc. A esa forma de recorrer se le llama \emph{Breadth-First
Search} (BFS, búsqueda en anchura), y para hacerla usamos una \emph{cola}: una fila de espera
donde el primero que entra es el primero que atendemos, igual que en la fila del
supermercado.

Lo único distinto aquí es que leemos cada piso en zigzag, como una serpiente:
el primer piso de izquierda a derecha, el siguiente de derecha a izquierda, el
siguiente otra vez de izquierda a derecha\ldots Truco para no perder tiempo: en
lugar de leer normal y luego voltear la lista, colocamos cada valor directamente
en el lugar que le toca dentro del piso. Si el piso tiene \texttt{capa} nodos, el
que sería el número \texttt{i} contando desde la izquierda lo guardamos en la
posición \texttt{capa-1-i} cuando toca ir al revés.

\begin{center}
\ttfamily
\begin{tabular}{c}
3\\
\end{tabular}
\quad
\begin{tabular}{cc}
9 & 20\\
\end{tabular}
\quad
\begin{tabular}{cc}
15 & 7\\
\end{tabular}
\end{center}

\noindent Árbol \texttt{[3,9,20,null,null,15,7]}: el nivel 0 es \texttt{[3]}
(izquierda a derecha), el nivel 1 se recorre de derecha a izquierda dando
\texttt{[20, 9]}, y el nivel 2 vuelve a izquierda a derecha con
\texttt{[15, 7]}. Resultado: \texttt{[[3], [20, 9], [15, 7]]}.

\paragraph{Cómo lo resuelve el código, paso a paso.} Recorremos el árbol por
niveles con una cola:
\begin{enumerate}
  \item Metemos la raíz en la cola.
  \item Mientras la cola no esté vacía, miramos cuántos nodos hay ahora mismo en
        ella: ese número es cuántos nodos tiene el piso actual.
  \item Sacamos justo esos nodos uno por uno, guardamos su valor en la lista del
        piso y metemos a la cola a sus hijos (que formarán el piso siguiente).
  \item Llevamos un interruptor de encendido/apagado (una \emph{bandera}) que
        cambia en cada piso. Cuando está encendido, guardamos los valores en la
        posición espejo para que ese piso quede al revés.
\end{enumerate}
Casos raros: si el árbol está vacío devolvemos una lista vacía de inmediato; si
solo hay un nodo, devolvemos ese único valor sin voltear nada, porque hay un
solo piso.

\lstinputlisting{src/02-arboles/103-zigzag-level-order.cpp}

\complejidad{$O(n)$}{$O(n)$}

\paragraph{¿Qué significa esa complejidad?} $n$ es el número total de nodos.
$O(n)$ en tiempo quiere decir que tocamos cada nodo una sola vez. $O(n)$ en
memoria es porque, en el peor caso, la cola y las listas del resultado llegan a
guardar todos los nodos a la vez. Es lo más rápido posible: no puedes escribir
todos los valores sin visitarlos al menos una vez.

\paragraph{Consejos para la entrevista (Amazon).} El error típico es leer el
piso normal y luego voltear la lista con una llamada aparte (\texttt{reverse});
funciona, pero colocar cada valor directo en su lugar ahorra esa vuelta extra y
demuestra que entiendes cómo está armado el piso. También puedes mencionar la
variante con dos pilas (una para cada dirección) como alternativa clásica.
Casos raros: árbol vacío, un solo nodo, y árboles muy torcidos donde algún piso
tiene un único nodo.
